\section{Ход работы}

\subsection{База данных библиотеки (пункт~\ref{task:db})}
Таблицы созданы скриптом \texttt{schema.sql} (листинг~\ref{lst:schema}). Связи между ними
показаны на рисунке~\ref{fig:er}.

\bbtCodeFile[sql, lines=3-23]{materials/schema.sql}{Создание таблиц}{lst:schema}

У каждой таблицы первичный ключ \texttt{id}. Таблица \texttt{loans} связывает читателя и
книгу внешними ключами \texttt{reader\_id} и \texttt{book\_id}, поэтому выдача
несуществующей книги или несуществующему читателю невозможна. Номер читательского билета
\texttt{ticket} уникален, год издания ограничен условием \texttt{CHECK}. Пустое поле
\texttt{returned} означает, что книга ещё на руках.

\bbtScreenshot[0.9\linewidth]{materials/er-diagram.png}{Схема базы данных
библиотеки}{fig:er}

Один читатель может иметь много выдач, и одна книга может выдаваться много раз: обе связи
«один ко многим» сходятся в таблице \texttt{loans}. В таблицы добавлены 4 читателя, 6 книг и
7 выдач. Выдачи, от которых зависят результаты запросов, приведены в
листинге~\ref{lst:loans}.

\bbtCodeFile[sql, lines=39-46]{materials/schema.sql}{Тестовые выдачи}{lst:loans}

\subsection{Условие и сортировка (пункт~\ref{task:where})}
Запрос в листинге~\ref{lst:q1} выбирает книги автора Миронова и сортирует их по году
издания по убыванию.

\bbtCodeFile[sql, lines=4-7]{materials/queries.sql}{Книги одного автора}{lst:q1}

Результат приведён в таблице~\ref{tab:q1}: у автора две книги, и книга 2018 года стоит
первой, как требует \texttt{ORDER BY year DESC}.

\begin{bbttable}{Результат запроса: книги автора Миронова}{tab:q1}
  \begin{tabular}{|l|c|}
    \hline
    title & year \\ \hline
    Основы SQL & 2018 \\ \hline
    Дискретная математика & 2017 \\ \hline
  \end{tabular}
\end{bbttable}

\subsection{Соединение таблиц (пункт~\ref{task:join})}
Запрос в листинге~\ref{lst:q2} соединяет выдачи с читателями и книгами и оставляет только
выдачи без даты возврата.

\bbtCodeFile[sql, lines=10-15]{materials/queries.sql}{Книги на руках}{lst:q2}

Из семи выдач возврата нет у четырёх, они и вошли в результат
(таблица~\ref{tab:q2}). Орлова Е. В. встречается дважды: у неё на руках две книги. Гусева П.
Р. в результате нет, потому что у него нет ни одной выдачи. Строки упорядочены по сроку
возврата, поэтому первым стоит самый ранний срок.

\begin{bbttable}{Результат запроса: книги на руках}{tab:q2}
  \begin{tabular}{|l|l|c|}
    \hline
    full\_name & title & due \\ \hline
    Белов К. Д. & Сети ЭВМ & 2026-10-13 \\ \hline
    Орлова Е. В. & Основы SQL & 2026-10-29 \\ \hline
    Савина М. А. & Операционные системы & 2026-11-17 \\ \hline
    Орлова Е. В. & Дискретная математика & 2026-11-18 \\ \hline
  \end{tabular}
\end{bbttable}

\subsection{Группировка (пункт~\ref{task:group})}
Запрос в листинге~\ref{lst:q3} считает выдачи каждого читателя и оставляет тех, у кого их
больше одной.

\bbtCodeFile[sql, lines=18-23]{materials/queries.sql}{Число выдач по читателям}{lst:q3}

Результат приведён в таблице~\ref{tab:q3}. Условие \texttt{HAVING} в этих данных ничего не
отбросило: читателя ровно с одной выдачей нет. Гусев П. Р. отсеялся раньше, при соединении
таблиц, поэтому строки с нулём выдач в результате нет; чтобы её получить, нужно
\texttt{LEFT JOIN} от таблицы читателей. Порядок Савиной М. А. и Белова К. Д. с одинаковым
числом выдач запрос не задаёт.

\begin{bbttable}{Результат запроса: число выдач}{tab:q3}
  \begin{tabular}{|l|c|}
    \hline
    full\_name & loans\_count \\ \hline
    Орлова Е. В. & 3 \\ \hline
    Савина М. А. & 2 \\ \hline
    Белов К. Д. & 2 \\ \hline
  \end{tabular}
\end{bbttable}

\subsection{Вычисление просрочки (пункт~\ref{task:late})}
Запрос в листинге~\ref{lst:q4} выбирает невозвращённые книги со сроком раньше даты проверки
и вычисляет число дней просрочки разностью \texttt{julianday}.

\bbtCodeFile[sql, lines=26-32]{materials/queries.sql}{Просроченные выдачи}{lst:q4}

Просрочены две выдачи из четырёх невозвращённых (таблица~\ref{tab:q4}). Книги Савиной М. А.
и вторая книга Орловой Е. В. в результат не попали: их сроки, 2026-11-17 и 2026-11-18, ещё не
наступили. Выдача Белову К. Д. книги «Основы SQL» тоже была просрочена, но книгу вернули,
и запрос её не учитывает.

\begin{bbttable}{Результат запроса: просроченные выдачи}{tab:q4}
  \begin{tabular}{|l|l|c|c|}
    \hline
    full\_name & title & due & days\_late \\ \hline
    Белов К. Д. & Сети ЭВМ & 2026-10-13 & 28 \\ \hline
    Орлова Е. В. & Основы SQL & 2026-10-29 & 12 \\ \hline
  \end{tabular}
\end{bbttable}

\subsection{Вложенный запрос (пункт~\ref{task:sub})}
Запрос в листинге~\ref{lst:q5} выбирает книги, номеров которых нет среди выданных.

\bbtCodeFile[sql, lines=35-38]{materials/queries.sql}{Книги, которые ни разу не
выдавались}{lst:q5}

Вложенный запрос возвращает номера всех выданных книг, а условие \texttt{NOT IN} оставляет
остальные. Ни разу не выдавалась только одна книга (таблица~\ref{tab:q5}). Столбец
\texttt{book\_id} объявлен как \texttt{NOT NULL}, поэтому \texttt{NOT IN} здесь работает
правильно: значение \texttt{NULL} во вложенном запросе сделало бы результат пустым.

\begin{bbttable}{Результат запроса: книги без выдач}{tab:q5}
  \begin{tabular}{|l|l|}
    \hline
    title & author \\ \hline
    Компиляторы & Карпов \\ \hline
  \end{tabular}
\end{bbttable}
